% Figure: the five principal classes of ECM macromolecule, with their
% mechanical role and tissue of greatest abundance.
\begin{figure}[htbp]
\centering
\begin{tikzpicture}[
  comp/.style={draw=engblack, rounded corners=2pt, minimum width=3.2cm,
    minimum height=1.0cm, align=center, font=\small},
  role/.style={font=\scriptsize, align=center, engblue},
]

\node[font=\bfseries] at (0, 3.2) {Principal components of the extracellular matrix};

\node[comp, fill=engred!15]   (col)   at (-5.5, 1.4) {Collagen\\ (I, II, III, IV)};
\node[comp, fill=englime!25]  (ela)   at (-2.7, 1.4) {Elastin};
\node[comp, fill=engblue!15]  (prot)  at ( 0.0, 1.4) {Proteoglycans\\ (aggrecan)};
\node[comp, fill=englime!15]  (gly)   at ( 2.7, 1.4) {Glycoproteins\\ (fibronectin)};
\node[comp, fill=engblue!10]  (ha)    at ( 5.5, 1.4) {Hyaluronic\\ acid};

\node[role] at (-5.5, 0.2) {tensile load,\\ $\sim 100\,\text{MPa}$};
\node[role] at (-2.7, 0.2) {reversible stretch,\\ $130\,\%$};
\node[role] at ( 0.0, 0.2) {water binding,\\ compression};
\node[role] at ( 2.7, 0.2) {cell-matrix\\ attachment};
\node[role] at ( 5.5, 0.2) {viscoelasticity,\\ lubrication};

\node[role, anchor=north] at (-5.5, -0.8) {bone, tendon, skin};
\node[role, anchor=north] at (-2.7, -0.8) {skin, lung, artery};
\node[role, anchor=north] at ( 0.0, -0.8) {cartilage};
\node[role, anchor=north] at ( 2.7, -0.8) {basement membrane};
\node[role, anchor=north] at ( 5.5, -0.8) {synovial fluid};

\end{tikzpicture}
\caption{Five principal classes of ECM macromolecule, the mechanical
role each carries, and the tissue in which each is dominant. The
engineer specifying a tissue scaffold matches not only the modulus of
the target tissue but its component-by-component matrix composition,
because each class signals to embedded cells through a different
receptor system.}
\label{fig:vol06ch05:ecm}
\end{figure}
